\subsection*{Reading and restoring paired verse}
A reader can prefer the original passage because of repeated glyphs or a change of literary style. We paired two passages within the same volume, with identical final two cells and different eight-glyph endings. The frozen image reader assigned both endings to the original passages on 1,276 of 2,131 joint likelihood comparisons (59.9\%, 95\% volume-bootstrap interval 56.4--63.9\%). Counting glyph overlap gave 58.5\% with half credit for ties. On the 347 pairs tied by overlap, the image reader preferred the original assignment on 195. The two-glyph reader tied on every pair. Requiring each passage to prefer its own ending independently reduced the count for the image reader to 208. Passage conditioning can therefore influence a comparison even when a frequent ending remains the more likely choice for both passages.

We next tested whether that information could guide multi-glyph writing. In 202,647 Tang verse pairs, masked second lines provided both a restoration task and training from a fully blank line. On 4,096 pairs from reserved poets, filling the most confident cell first recovered 761 of 23,114 missing glyphs from a blank line. Left-first filling with the same weights recovered 744. The difference of 0.07 percentage points had an author-bootstrap interval from $-0.05$ to 0.16. Each method completed zero reference lines. The byte writer recovered 916 glyphs and aligned counts recovered 900. With two adjacent cells missing, the visual writer completed seven lines, all of which failed after exchanging the first line.

The image encoder had received a weak learning signal for the identity of each visible glyph. We therefore trained it to recover the canonical bitmap of every source cell and every surviving target cell, using the same classifier as for writing. Across 6,362 lines from 156 separately reserved poet attributions, this change raised blank-line recovery from 1,281 to 1,392 of 37,442 glyphs (gain 0.30 percentage points, 95\% interval 0.17--0.42). The matched byte writer recovered 1,483 and aligned counts recovered 1,587. None recovered a complete blank reference. Two-cell repairs rose from 456 to 502 correct glyphs, while complete repairs fell from 20 to 17. Figure~\ref{fig:verse-canvas} shows an exact two-cell repair and the repeated phrasing of a blank-line prediction.

With the first line in an unseen typeface, nearest-neighbour identification in the learned source representation rose from 1,651 to 5,519 of 8,414 glyphs. Raw ink-density features identified 5,475. Despite the improvement in recognition, both visual writers recovered exactly 857 blank-line glyphs in that typeface. A description of visual learning must distinguish preserving written identity from choosing a suitable next expression.
